\documentclass{article}
\usepackage{amsmath,amssymb}
\usepackage{xurl}
\usepackage[hidelinks]{hyperref}
% Shared commands for notes in docs/paper-gaps/.

\DeclareUnicodeCharacter{2080}{\ifmmode{}_{0}\else\textsubscript{0}\fi}
\DeclareUnicodeCharacter{2081}{\ifmmode{}_{1}\else\textsubscript{1}\fi}
\DeclareUnicodeCharacter{2082}{\ifmmode{}_{2}\else\textsubscript{2}\fi}
\DeclareUnicodeCharacter{2099}{\ifmmode{}_{n}\else\textsubscript{n}\fi}

\newcommand{\Lean}{\textsc{Lean}}
\ExplSyntaxOn
\NewDocumentCommand{\leanid}{m}{
  \ifmmode
    \textsf{\detokenize{#1}}
  \else
    \group_begin:
    \urlstyle{sf}
    \tl_set:Nn \l_tmpa_tl {#1}
    \tl_replace_all:Nnn \l_tmpa_tl {\_} {_}
    \exp_args:NV \path \l_tmpa_tl
    \group_end:
  \fi
}
\ExplSyntaxOff
\providecommand{\tr}{\operatorname{tr}}
\newcommand{\tnleanIssueBase}{https://github.com/LionSR/TNLean/issues/}
\newcommand{\tnleanPrBase}{https://github.com/LionSR/TNLean/pull/}
\newcommand{\tnleanDocBase}{https://sirui-lu.com/TNLean/docs/}
\newcommand{\ghissue}[1]{\href{\tnleanIssueBase#1}{GitHub issue \##1}}
\newcommand{\ghpr}[1]{\href{\tnleanPrBase#1}{PR \##1}}
\newcommand{\doclink}[1]{\href{\tnleanDocBase#1.html}{\path{#1}}}

% Dirac notation and the identity operator, as in blueprint/src/macros/common.tex.
% \providecommand keeps note-local definitions authoritative.
\usepackage{dsfont}
\providecommand{\ket}[1]{|#1\rangle}
\providecommand{\bra}[1]{\langle #1|}
\providecommand{\braket}[2]{\langle #1 | #2 \rangle}
\providecommand{\ketbra}[2]{|#1\rangle\!\langle #2|}
\providecommand{\Id}{\mathds{1}}
\providecommand{\one}{\mathds{1}}
\providecommand{\C}{\mathbb{C}}
\providecommand{\MN}[1]{M_{#1}(\C)}
\providecommand{\MD}{M_D(\C)}

% Verdict marker: \gapnote{<kind>}{<status>}. Kind is the policy
% classification (clarification, local-correction, scope-restriction,
% unfaithful, false-source, open-gap); status is open, wip (elimination
% underway), resolved, or historical. Machine-read by the site index;
% prints a verdict line.
\newcommand{\gapnote}[2]{\par\noindent\textbf{Verdict:} #1 (#2).\par}

\title{Virtual Bond Dimension in MPU Tensor Equivalence}
\author{}
\date{2026-08-10, revised 2026-10-06}
\begin{document}
\maketitle
\gapnote{local-correction}{resolved}

\section{Source definition}

Let $\mathcal U$ and $\mathcal V$ be matrix product unitary tensors with
physical dimensions $d_a$ and $d_b$. Let $p_a$ and $p_b$ be positive coprime
ancilla dimensions. Cirac--P\'erez-Garc\'ia--Schuch--Verstraete define strict
equivalence of two tensors in canonical form through a continuous path of
tensors from one to the other, not necessarily in canonical form. They define
equivalence by first attaching identity ancillas and then applying a common
blocking length \cite[lines~706--724]{Cirac2017MPU}. The enlarged physical
dimensions satisfy
\begin{align}
    p_a d_a &= p_b d_b.
    \label{eq:mpu_equivalence_fixed_bond_physical_dimension_balance}
\end{align}
The cited definition does not say in which space the path runs when the two
tensors have different virtual bond dimensions.

\section{Why the source needs no comment and the library does}

The canonical form of the source (its Section~II) admits blocks of weight
zero. A tensor with unused bond directions adjoined, its letters placed in the
upper left corner of a larger bond space and zero elsewhere, is still in the
source's canonical form, so two tensors of different bond dimensions are
compared inside the larger bond space without further remark. The library's
canonical form (\leanid{MPSTensor.IsMPUCanonicalForm}) requires nonzero
weights and blocks filling the bond space, the nonzero-coefficient convention
of \texttt{docs/audits/2026-08-23\_nonzero\_coefficient\_convention.md}. Under
that convention a tensor with adjoined directions is not in canonical form,
and a definition that requires canonical form of its endpoints has to say that
the path joins the enlarged tensors rather than the tensors themselves.

\section{Formal convention}

For $D\le D'$, \leanid{MPOTensor.padBond} adjoins unused bond directions to a
tensor of bond dimension $D$ up to bond dimension $D'$: every letter is the
original letter in the upper left $D\times D$ corner and zero elsewhere. By
\leanid{MPOTensor.mpo\_padBond} this changes no periodic operator of positive
length, so it preserves the MPU property (\leanid{MPOTensor.IsMPU.padBond})
and the index (\leanid{MPOTensor.IsMPU.index\_padBond}).

The declarations \leanid{MPOTensor.StrictlyEquivalent} and
\leanid{MPOTensor.StrictlyEquivalentUnderSymmetry} require canonical form of
the two given tensors, of bond dimensions $D_a$ and $D_b$, and a continuous
path of MPU tensors of some bond dimension $D'\ge\max(D_a,D_b)$ from the
enlargement of the first to the enlargement of the second. The declarations
\leanid{MPOTensor.Equivalent} and \leanid{MPOTensor.EquivalentUnderSymmetry}
apply this after ancillas and blocking as the source prescribes; canonical
form of the blocked tensors with ancilla is a condition, as in the source.
The former definitions, which fixed one bond dimension $D$ at the endpoints
and along the path, are the case $D'=D_a=D_b$
(\leanid{MPOTensor.StrictlyEquivalent.of\_fixedBond},
\leanid{MPOTensor.StrictlyEquivalentUnderSymmetry.of\_fixedBond}).

The placement of the letters in the upper left corner is a choice. Any other
placement differs from it by a unitary of the larger bond space, and the
unitary group is connected, so the relation does not depend on the placement.

\section{Verdict}

\textbf{Verdict: local correction, resolved.} The source's silence on the
bond space of the path is filled by the convention that two tensors are
compared after unused bond directions are adjoined to each up to a common
bond dimension. This is the reading under which the source's own canonical
form makes its definition meaningful, transported to the library's canonical
form. No other stabilization is introduced. The earlier scope restriction
(one fixed bond dimension at the endpoints and along the path) is removed.

\bibliographystyle{alpha}
\bibliography{references}
\end{document}
